% GENERATED FILE. Edit canonical lessons, then run npm run generate:books.
% canonical-source-hash: fnv1a64:5b9d38761b9694f8
% canonical-lessons: LA-C395-maturo, LA-C395-differo, LA-C395-prorogo, LA-C395-dilatio, LA-C395-exspectatio

\chapter{Sooner or Later}
\label{ch:la-sooner-or-later}
\hlchaptermodality{395}

\begin{chapteropening}
\textbf{By the end of this chapter:} \emph{I can say that I hurry something on or put it off, that a command was extended, and talk about a delay and what people expect.}

Complete the last lesson of chapter 395: I can say that I hurry something on or put it off, that a command was extended, and talk about a delay and what people expect.
\end{chapteropening}

\section[mātūrō, mātūrāre]{mātūrō, mātūrāre — to hasten; to ripen}

\label{lesson:LA-C395-maturo}



\begin{quote}
Before the new one: say the Latin for to set about, then the Latin for to foresee.
\end{quote}

\subsection*{You'll want to know: mātūrō, mātūrāre}
\textbf{mātūrō, mātūrāre} — \textquotedblleft{}to hasten; to ripen\textquotedblright{}.

\begin{quote}
Iter mātūrāvit.
\end{quote}

\begin{quote}
\emph{He hurried on his journey.}
\end{quote}

\emph{mātūrō, mātūrāre, mātūrāvī, mātūrātum}, built on \emph{mātūrus}, ripe, early: to bring on early. English \emph{mature}.

\begin{cousinweb}
\emph{mātūrō} is built on \emph{mātūrus} \textquotedblleft{}ripe, early, timely\textquotedblright{} (EtymDB; Lewis \& Short). English \emph{mature} comes from the adjective.
\end{cousinweb}

\begin{grammarlens}[title={What you've built}]
One more word for telling, explaining and describing at B1.
\end{grammarlens}

\subsection*{Your turn}
\begin{itemize}
\raggedright
  \item \emph{Say it:} \emph{mātūrō, mātūrāre}
  \item \emph{Say it:} \emph{mātūrō, mātūrāre}, once more
  \item \emph{Say it:} a sentence of your own with \emph{mātūrō, mātūrāre}
  \item \emph{Recall:} say the Latin for to set about, then the Latin for to foresee, then say \emph{mātūrō, mātūrāre} again
\end{itemize}

\begin{tcolorbox}[breakable,colback=teal!4,colframe=teal!35!black,title={Before you move on}]
What does mātūrō, mātūrāre mean? (\textquotedblleft{}To hasten; to ripen\textquotedblright{}.) Say it once more.
\end{tcolorbox}
\section[differō, differre]{differō, differre — to put off; to differ; to spread}

\label{lesson:LA-C395-differo}



\begin{quote}
Before the new one: say the Latin for to foresee, then the Latin for to hasten.
\end{quote}

\subsection*{You'll want to know: differō, differre}
\textbf{differō, differre} — \textquotedblleft{}to put off; to differ; to spread\textquotedblright{}.

\begin{quote}
Rem in crāstinum diem distulimus.
\end{quote}

\begin{quote}
\emph{We put the matter off to the next day.}
\end{quote}

\emph{differō, differre, distulī, dīlātum}: irregular, like \emph{ferō}. To put off, and, without an object, to differ. English \emph{defer} and \emph{differ}.

\begin{cousinweb}
\emph{differō} is built from \emph{dis-} \textquotedblleft{}apart\textquotedblright{} and \emph{ferō} \textquotedblleft{}to carry\textquotedblright{} (EtymDB; Lewis \& Short): to carry in different directions, or forward in time. English \emph{defer} and \emph{differ} both come from it.
\end{cousinweb}

\begin{grammarlens}[title={What you've built}]
One more word for telling, explaining and describing at B1.
\end{grammarlens}

\subsection*{Your turn}
\begin{itemize}
\raggedright
  \item \emph{Say it:} \emph{differō, differre}
  \item \emph{Say it:} \emph{differō, differre}, once more
  \item \emph{Say it:} a sentence of your own with \emph{differō, differre}
  \item \emph{Recall:} say the Latin for to foresee, then the Latin for to hasten, then say \emph{differō, differre} again
\end{itemize}

\begin{tcolorbox}[breakable,colback=teal!4,colframe=teal!35!black,title={Before you move on}]
What does differō, differre mean? (\textquotedblleft{}To put off; to differ; to spread\textquotedblright{}.) Say it once more.
\end{tcolorbox}
\section[prōrogō, prōrogāre]{prōrogō, prōrogāre — to extend, to prolong (a term of office)}

\label{lesson:LA-C395-prorogo}



\begin{quote}
Before the new one: say the Latin for to hasten, then the Latin for to put off.
\end{quote}

\subsection*{You'll want to know: prōrogō, prōrogāre}
\textbf{prōrogō, prōrogāre} — \textquotedblleft{}to extend, to prolong (a term of office)\textquotedblright{}.

\begin{quote}
Senātus imperium eī prōrogāvit.
\end{quote}

\begin{quote}
\emph{The senate extended his command.}
\end{quote}

\emph{prōrogō, prōrogāre, prōrogāvī, prōrogātum}, built on \emph{rogō}: to ask the people for more time. English \emph{prorogue}.

\begin{cousinweb}
\emph{prōrogō} is built from \emph{prō} \textquotedblleft{}forward\textquotedblright{} and \emph{rogō} \textquotedblleft{}to ask\textquotedblright{} (EtymDB; Lewis \& Short): to ask for a term to run on. English \emph{prorogue} comes from it.
\end{cousinweb}

\begin{grammarlens}[title={What you've built}]
One more word for telling, explaining and describing at B1.
\end{grammarlens}

\subsection*{Your turn}
\begin{itemize}
\raggedright
  \item \emph{Say it:} \emph{prōrogō, prōrogāre}
  \item \emph{Say it:} \emph{prōrogō, prōrogāre}, once more
  \item \emph{Say it:} a sentence of your own with \emph{prōrogō, prōrogāre}
  \item \emph{Recall:} say the Latin for to hasten, then the Latin for to put off, then say \emph{prōrogō, prōrogāre} again
\end{itemize}

\begin{tcolorbox}[breakable,colback=teal!4,colframe=teal!35!black,title={Before you move on}]
What does prōrogō, prōrogāre mean? (\textquotedblleft{}To extend, to prolong (a term of office)\textquotedblright{}.) Say it once more.
\end{tcolorbox}
\section[dīlātiō, dīlātiōnis]{dīlātiō, dīlātiōnis — a putting off, a delay}

\label{lesson:LA-C395-dilatio}



\begin{quote}
Before the new one: say the Latin for to put off, then the Latin for to extend.
\end{quote}

\subsection*{You'll want to know: dīlātiō, dīlātiōnis}
\textbf{dīlātiō, dīlātiōnis} — \textquotedblleft{}a putting off, a delay\textquotedblright{}.

\begin{quote}
Rēs dīlātiōnem nōn patitur.
\end{quote}

\begin{quote}
\emph{The matter allows no delay.}
\end{quote}

\emph{dīlātiō, dīlātiōnis} f., the noun of \emph{differō}, from its supine \emph{dīlātum}. English \emph{dilatory}.

\begin{cousinweb}
\emph{dīlātiō} is the noun made on \emph{dīlātum}, the supine of \emph{differō} \textquotedblleft{}to put off\textquotedblright{} (Lewis \& Short), whose \emph{-lātum} comes from a different old verb than \emph{ferō}.
\end{cousinweb}

\begin{grammarlens}[title={What you've built}]
One more word for telling, explaining and describing at B1.
\end{grammarlens}

\subsection*{Your turn}
\begin{itemize}
\raggedright
  \item \emph{Say it:} \emph{dīlātiō, dīlātiōnis}
  \item \emph{Say it:} \emph{dīlātiō, dīlātiōnis}, once more
  \item \emph{Say it:} a sentence of your own with \emph{dīlātiō, dīlātiōnis}
  \item \emph{Recall:} say the Latin for to put off, then the Latin for to extend, then say \emph{dīlātiō, dīlātiōnis} again
\end{itemize}

\begin{tcolorbox}[breakable,colback=teal!4,colframe=teal!35!black,title={Before you move on}]
What does dīlātiō, dīlātiōnis mean? (\textquotedblleft{}A putting off, a delay\textquotedblright{}.) Say it once more.
\end{tcolorbox}
\section[exspectātiō, exspectātiōnis]{exspectātiō, exspectātiōnis — expectation, anticipation}

\label{lesson:LA-C395-exspectatio}



\begin{quote}
Before the new one: say the Latin for to extend, then the Latin for a putting off.
\end{quote}

\subsection*{You'll want to know: exspectātiō, exspectātiōnis}
\textbf{exspectātiō, exspectātiōnis} — \textquotedblleft{}expectation, anticipation\textquotedblright{}.

\begin{quote}
Omnēs in magnā exspectātiōne erant.
\end{quote}

\begin{quote}
\emph{Everyone was in great expectation.}
\end{quote}

\emph{exspectātiō, exspectātiōnis} f., the noun of \emph{exspectō}. \emph{Praeter exspectātiōnem}, beyond expectation. English \emph{expectation}.

\begin{cousinweb}
\emph{exspectātiō} is the noun built on \emph{exspectō} \textquotedblleft{}to wait for\textquotedblright{} (Lewis \& Short), from \emph{ex} \textquotedblleft{}out\textquotedblright{} and \emph{spectō} \textquotedblleft{}to look\textquotedblright{}: to look out for. English \emph{expectation} comes from it.
\end{cousinweb}

\begin{grammarlens}[title={What you've built}]
One more word for telling, explaining and describing at B1.
\end{grammarlens}

\subsection*{Your turn}
\begin{itemize}
\raggedright
  \item \emph{Say it:} \emph{exspectātiō, exspectātiōnis}
  \item \emph{Say it:} \emph{exspectātiō, exspectātiōnis}, once more
  \item \emph{Say it:} a sentence of your own with \emph{exspectātiō, exspectātiōnis}
  \item \emph{Recall:} say the Latin for to extend, then the Latin for a putting off, then say \emph{exspectātiō, exspectātiōnis} again
\end{itemize}

\begin{tcolorbox}[breakable,colback=teal!4,colframe=teal!35!black,title={Before you move on}]
What does exspectātiō, exspectātiōnis mean? (\textquotedblleft{}Expectation, anticipation\textquotedblright{}.) Say it once more.
\end{tcolorbox}
